\chapter{Permukaan}\label{ch:b2:surfaces}

Setelah kurva, kini permukaan: yakni objek berparameter dua di $\R^3$. Adapun
kalkulus \omterm{def:b2:diffcalc:differential}{diferensial} pada \cref{ch:b2:diffcalc} menyediakan segala yang
kita perlukan --- sebab turunan parsial memberi vektor singgung, hasil kali
silang memberi normalnya, dan \omterm{def:b2:linalg:det}{determinan} memberi luasnya. Kita mendefinisikan
\omterm{def:b2:surfaces:param}{permukaan terparameter} yang regular, \omterm{def:b2:surfaces:tangent}{bidang singgungnya}, dan
\emph{\omterm{def:b2:surfaces:fff}{bentuk fundamental pertama}}, yang menyandikan semua pengukuran \omterm{def:b2:curves:length}{panjang}
dan \omterm{def:b2:surfaces:area}{luas} pada permukaannya. Permukaan juga muncul sebagai himpunan aras
$f(x, y, z) = c$; dan gradiennya lalu mengarahkan normalnya.

\section{Permukaan terparameter}

\begin{definition}[Permukaan terparameter yang regular]\label{def:b2:surfaces:param}
Misalkan $U \subseteq \R^2$ \omterm{def:b2:metric:topology}{terbuka}. Sebuah \emph{permukaan
terparameter}\index{permukaan terparameter} berkelas
$\mathcal{C}^k$ (dengan $k \geq 1$) adalah pemetaan $\sigma \colon U \to
\R^3$, $(u, v) \mapsto \sigma(u, v)$, yang berkelas $\mathcal{C}^k$. Sebuah
titik disebut \emph{regular} bila vektor turunan parsialnya
\[
\sigma_u(u, v) = \frac{\partial\sigma}{\partial u}(u,v),
\qquad
\sigma_v(u, v) = \frac{\partial\sigma}{\partial v}(u,v)
\]
bebas linear, yakni $\sigma_u \wedge \sigma_v \neq 0$;
dan permukaannya disebut \emph{regular} bila setiap titiknya demikian.
\end{definition}

\begin{example}[Tiga pemerian yang baku]\label{ex:b2:surfaces:standard}
\leavevmode
\begin{enumerate}
\item \emph{Grafik:} $\sigma(u, v) = (u,\ v,\ f(u, v))$ untuk $f \in
\mathcal{C}^1(U)$. Selalu regular: sebab $\sigma_u = (1, 0, f_u)$ dan
$\sigma_v = (0, 1, f_v)$ bersifat bebas.
\item \emph{Bola (dalam koordinat bola):} untuk bola berjari-jari
$R$,
\[
\sigma(\theta, \varphi)
= (R\cos\theta\cos\varphi,\ R\sin\theta\cos\varphi,\
R\sin\varphi),
\qquad
(\theta, \varphi) \in \R \times
\bigl(-\tfrac\pi2, \tfrac\pi2\bigr),
\]
dengan $\theta$ menyatakan bujurnya dan $\varphi$ lintangnya. Kita memeriksa
$\norm{\sigma_\theta \wedge \sigma_\varphi} = R^2\cos\varphi > 0$:
jadi regular jauh dari kutubnya (yang dihilangkan bagan ini).
\item \emph{Himpunan aras:} $S = \{(x,y,z) : f(x, y, z) = c\}$ dengan
$f$ bersifat $\mathcal{C}^1$ dan $\nabla f \neq 0$ pada $S$. Di dekat tiap
titiknya, satu koordinatnya dapat dinyatakan sebagai fungsi dua
lainnya lewat teorema fungsi implisit (\cref{ch:b2:diffcalc}), jadi
$S$ secara lokal berupa grafik.
\end{enumerate}
\end{example}

\begin{example}[Dari himpunan aras ke grafik]
\label{ex:b2:surfaces:levelgraph}
Teorema fungsi implisit pada butir 3 pantas dijalankan sekali
secara gamblang. Ambillah bola $x^2 + y^2 + z^2 = R^2$ di dekat kutub
utaranya $(0, 0, R)$: di sana $\frac{\partial f}{\partial z} = 2z =
2R \neq 0$, jadi menyelesaikannya untuk $z$ memberi bagan grafiknya
\[
z = \sqrt{R^2 - x^2 - y^2},
\qquad x^2 + y^2 < R^2 ,
\]
yang regular di mana-mana pada daerah (\omterm{def:b2:metric:topology}{terbuka}) definisinya --- termasuk
di kutub yang dilewatkan bagan bolanya. Adapun di dekat titik khatulistiwa
seperti $(R, 0, 0)$ teorema yang sama menyelesaikannya untuk $x$
(sebab $\frac{\partial f}{\partial x} = 2R \neq 0$). Jadi patokan
praktisnya: sebuah permukaan aras berupa grafik di atas bidang koordinat
yang \emph{ortogonal terhadap komponen gradien terbesarnya},
dan dengan meliputi bolanya lewat enam bagan grafik semacam itu kita
memeriksa kemulusannya di mana-mana tanpa trigonometri sama
sekali.
\end{example}

\section{Bidang singgung dan normalnya}

\begin{definition}[Bidang singgung]\label{def:b2:surfaces:tangent}
Misalkan $\sigma$ regular di $(u_0, v_0)$, dengan $M_0 = \sigma(u_0, v_0)$.
Adapun \emph{bidang singgung}\index{bidang singgung} $T_{M_0}S$ adalah bidang lewat $M_0$
yang diarahkan $\operatorname{Vect}(\sigma_u, \sigma_v)$ (dengan turunan parsialnya di
$(u_0, v_0)$). Sedangkan \emph{normal satuannya}\index{normal satuan} adalah
\[
n(u_0, v_0)
= \frac{\sigma_u \wedge \sigma_v}
       {\norm{\sigma_u \wedge \sigma_v}} .
\]
\end{definition}

\begin{proposition}[Vektor singgung itu vektor
kecepatan]\label{prop:b2:surfaces:velocity}
Arah $T_{M_0}S$ persis merupakan himpunan vektor
$\gamma'(0)$, dengan $\gamma = \sigma \circ c$ menjelajahi
kurva $\mathcal{C}^1$ yang digambar pada permukaannya lewat $M_0$ (yakni
$c \colon (-\varepsilon, \varepsilon) \to U$ bersifat $\mathcal{C}^1$
dengan $c(0) = (u_0, v_0)$).
\end{proposition}

\begin{proof}
Bila $c(t) = (u(t), v(t))$, maka aturan rantai
(\cref{ch:b2:diffcalc}) memberi
\[
\gamma'(0) = u'(0)\,\sigma_u + v'(0)\,\sigma_v
\in \operatorname{Vect}(\sigma_u, \sigma_v).
\]
Sebaliknya, vektor $a\sigma_u + b\sigma_v$ dicapai oleh
kurva $c(t) = (u_0 + at,\ v_0 + bt)$, yang tinggal di himpunan \omterm{def:b2:metric:topology}{terbuka}
$U$ untuk $\abs t$ yang kecil.
\end{proof}

\begin{example}[Bidang singgung helikoidnya]
\label{ex:b2:surfaces:helicoidtangent}
Untuk helikoid $\sigma(u, v) = (v\cos u,\ v\sin u,\ au)$,
di titik $\sigma(0, 1) = (1, 0, 0)$:
\[
\sigma_u = (0,\ 1,\ a), \qquad \sigma_v = (1,\ 0,\ 0),
\qquad
\sigma_u \wedge \sigma_v = (0,\ a,\ -1),
\]
jadi \omterm{def:b2:surfaces:tangent}{bidang singgungnya} adalah $a\,y = z$. Ia memuat seluruh
garis pembangun mendatarnya $\{(t, 0, 0)\}$ (yakni arah $\sigma_v$): sebab
seperti pada kerucut \cref{exo:b2:surfaces:1}, permukaan yang \omterm{def:b2:structures:generated}{dibangun}
garis lurus punya tiap garis pembangunnya terletak di dalam \omterm{def:b2:surfaces:tangent}{bidang
singgung} sepanjangnya. Adapun arah singgung lainnya $\sigma_u$ adalah
kecepatan heliks $u \mapsto \sigma(u, 1)$: jadi satu bagan,
dua kurva tergambar, dan seluruh \omterm{def:b2:surfaces:tangent}{bidang singgungnya} terentang ---
yakni \cref{prop:b2:surfaces:velocity} yang beraksi.
\end{example}

\begin{proposition}[Normal sebuah permukaan aras]\label{prop:b2:surfaces:gradient}
Misalkan $S = \{f = c\}$ dengan $f$ berkelas $\mathcal{C}^1$ dan $\nabla
f(M_0) \neq 0$. Maka \omterm{def:b2:surfaces:tangent}{bidang singgung} $S$ di $M_0$ adalah
bidang lewat $M_0$ yang ortogonal terhadap $\nabla f(M_0)$:
\[
T_{M_0}S :\quad
\langle \nabla f(M_0),\ M - M_0\rangle = 0 .
\]
\end{proposition}

\begin{proof}
Untuk sembarang kurva $\gamma$ yang digambar pada $S$ lewat $M_0$, berlaku $f(\gamma(t)) =
c$ secara identik, jadi aturan rantainya memberi $\langle
\nabla f(M_0), \gamma'(0)\rangle = 0$: sehingga semua vektor kecepatannya
ortogonal terhadap gradiennya, jadi arah singgungnya termuat
di bidang $\nabla f(M_0)^\perp$. Keduanya subruang
berdimensi $2$ --- yakni arah singgungnya karena $S$ secara lokal berupa
grafik yang regular (\cref{ex:b2:surfaces:standard}), dan komplemen
ortogonalnya karena $\nabla f(M_0) \neq 0$ --- sehingga keduanya
sama.
\end{proof}

\begin{example}\label{ex:b2:surfaces:spheretangent}
Untuk bola $x^2 + y^2 + z^2 = R^2$: berlaku $\nabla f = 2(x, y, z)$,
jadi \omterm{def:b2:surfaces:tangent}{bidang singgung} di $M_0$ ortogonal terhadap jari-jari
$\vect{OM_0}$ --- yakni fakta klasik bahwa jari-jari dan \omterm{def:b2:surfaces:tangent}{bidang singgungnya}
saling tegak lurus, dengan persamaan $\langle M_0, M\rangle = R^2$.
\end{example}

\begin{example}[Bidang singgung sebuah grafik]\label{ex:b2:surfaces:graphtangent}
Untuk $z = f(x, y)$ di $(x_0, y_0)$: dengan menerapkan
\cref{prop:b2:surfaces:gradient} pada $F(x,y,z) = f(x,y) - z$,
\[
z = f(x_0, y_0)
+ f_x(x_0, y_0)(x - x_0)
+ f_y(x_0, y_0)(y - y_0) ,
\]
yakni bagian \omterm{def:b2:affine:subspace}{afin} uraian Taylor orde pertamanya --- jadi
\omterm{def:b2:surfaces:tangent}{bidang singgungnya} adalah grafik \omterm{def:b2:diffcalc:differential}{diferensialnya}, sebagaimana mestinya.
\end{example}

\begin{example}[Titik terdekat sebuah
permukaan]\label{ex:b2:surfaces:closest}
Titik manakah pada paraboloid $z = x^2 + y^2$ yang terdekat ke
$P = (0, 0, 1)$? Minimumkanlah kuadrat jaraknya sepanjang
permukaannya: dengan $\rho^2 = x^2 + y^2$,
\[
g(\rho^2) = \rho^2 + (\rho^2 - 1)^2,
\qquad
g'(\rho^2) = 1 + 2(\rho^2 - 1) = 0
\iff \rho^2 = \tfrac12 ,
\]
yang memberi lingkaran titik pada ketinggian $z = \frac12$ dan
berjarak $\sqrt{\tfrac12 + \tfrac14} = \frac{\sqrt3}2$. Adapun
tanda geometris keminimalannya: di titik $M$ semacam itu,
vektor $\vect{MP}$ mestilah \emph{normal} terhadap permukaannya ---
sebab kalau tidak, meluncur sepanjang kurva tergambar yang kecepatannya punya
komponen ke arah $P$ akan mengecilkan jaraknya. Periksa:
$\nabla(z - x^2 - y^2) = (-2x, -2y, 1)$ di $M = (x, y,
\tfrac12)$, sedangkan $\vect{MP} = (-x, -y, \tfrac12) =
\tfrac12(-2x, -2y, 1)$: jadi sejajar, seperti diramalkan. Adapun
syarat orde pertama ``kaki garis tegak lurusnya'' sama dengan
yang kelak menggerakkan ekstremum pada himpunan aras di
\cref{exo:b2:surfaces:12}.
\end{example}

\begin{example}[Bidang singgung kuadrik: aturan
polarisasi]\label{ex:b2:surfaces:polarization}
Misalkan $S : xy + yz + zx = 1$ dan $M_0 = (1, 1, 0) \in S$. Di sini
$\nabla f = (y + z,\ x + z,\ x + y)$, jadi $\nabla f(M_0) = (1,
1, 2)$ dan \omterm{def:b2:surfaces:tangent}{bidang singgungnya} adalah
\[
(x - 1) + (y - 1) + 2z = 0,
\qquad\text{yakni}\qquad x + y + 2z = 2 .
\]
Jawaban yang sama datang dari \emph{aturan \omterm{def:b2:quadratic:def}{polarisasi}} yang
menyamaratakan \cref{ex:b2:surfaces:spheretangent} dan
\cref{exo:b2:surfaces:2}: pada persamaan \omterm{pb:b2:surfaces:1}{kuadriknya}, gantilah
$x^2$ dengan $x_0x$, dan tiap hasil kali $xy$ dengan $\frac{x_0y +
y_0x}2$ (dan secara \omterm{def:b2:structures:generated}{siklik}):
\[
\frac{x_0y + y_0x}2 + \frac{y_0z + z_0y}2 + \frac{z_0x +
x_0z}2 = 1
\;\xrightarrow{\ M_0 = (1,1,0)\ }\;
\frac{x + y}2 + \frac z2 + \frac z2 = 1 ,
\]
yang kembali menjadi $x + y + 2z = 2$. Aturannya berhasil karena
$\nabla$ sebuah \omterm{def:b2:quadratic:def}{bentuk kuadratik} adalah bentuk bilinear yang bersesuaian
yang dinilai terhadap titik pangkalnya --- jadi kesinggungan pada \omterm{pb:b2:surfaces:1}{kuadrik} adalah
\omterm{def:b2:quadratic:def}{polarisasi}, yakni satu wajah lagi \cref{ch:b2:quadratic}.
\end{example}

\section{Bentuk fundamental pertama}

\begin{definition}[Bentuk fundamental pertama]\label{def:b2:surfaces:fff}
Misalkan $\sigma \colon U \to \R^3$ permukaan $\mathcal{C}^1$ yang
regular. Adapun \emph{bentuk fundamental pertamanya}\index{bentuk fundamental
pertama} di $(u,v)$ adalah
\omterm{def:b2:quadratic:def}{bentuk kuadratik} definit positif pada $\R^2$
\[
I(h, k) = \norm{h\,\sigma_u + k\,\sigma_v}^2
= E\,h^2 + 2F\,hk + G\,k^2,
\]
dengan
\[
E = \norm{\sigma_u}^2,
\qquad
F = \langle \sigma_u, \sigma_v\rangle,
\qquad
G = \norm{\sigma_v}^2 .
\]
\end{definition}

\begin{remark}
Bentuk $I$ adalah pembatasan hasil kali skalar Euklides ruang lingkupnya ke
\omterm{def:b2:surfaces:tangent}{bidang singgungnya}, yang dibaca dalam basis $(\sigma_u, \sigma_v)$: jadi ia
definit positif persis karena $\sigma_u, \sigma_v$ bersifat
bebas (\cref{ch:b2:quadratic}). Adapun setiap besaran metrik pada
permukaannya --- yakni \omterm{def:b2:curves:length}{panjang} kurva tergambar, sudut antar keduanya, \omterm{def:b2:surfaces:area}{luas} ---
dihitung dari $E, F, G$ belaka. Jadi dua permukaan dengan $E,
F, G$ yang sama dalam parameter yang sesuai bersifat \emph{isometrik} walaupun keduanya duduk
berbeda di ruangnya: dan inilah titik tolak geometri
yang hakiki.
\end{remark}

\begin{example}[Sudut antar kurva
koordinat]\label{ex:b2:surfaces:coordangle}
\omterm{def:b2:surfaces:fff}{Bentuk fundamental pertamanya} juga mengukur sudut: sebab
kurva koordinat $u \mapsto \sigma(u, v_0)$ dan $v \mapsto
\sigma(u_0, v)$ bertemu pada sudut $\theta$ dengan
\[
\cos\theta = \frac{\langle\sigma_u,
\sigma_v\rangle}{\norm{\sigma_u}\,\norm{\sigma_v}}
= \frac{F}{\sqrt{EG}} :
\]
jadi satu koefisien $F$ belaka yang memutuskan keortogonalan
jaring parameternya. Untuk bagan bola dan helikoidnya, $F =
0$: sehingga meridiannya memotong paralelnya, dan heliksnya memotong garis
pembangun mendatarnya, secara siku-siku --- dan itulah sebabnya integran
luasnya runtuh menjadi $\sqrt{EG}$. Sedangkan untuk bagan grafik,
$F = f_xf_y$ hanya lenyap di tempat sebuah turunan parsialnya lenyap:
jadi jaring koordinat sebuah grafik yang miring \emph{tak}
ortogonal, walaupun jaring $(x, y)$ di bawahnya ortogonal. Jadi ketika
perhitungan pada sebuah permukaan tampak berat, langkah pertamanya adalah
mencari bagan dengan $F = 0$.
\end{example}

\begin{example}[Bagan pelananya]\label{ex:b2:surfaces:saddlechart}
Untuk pelana $z = xy$ dengan bagan $\sigma(u, v) = (u, v,
uv)$:
\[
\sigma_u = (1, 0, v), \qquad \sigma_v = (0, 1, u),
\qquad
E = 1 + v^2, \quad F = uv, \quad G = 1 + u^2 ,
\]
dan $EG - F^2 = 1 + u^2 + v^2 > 0$: jadi regular di mana-mana. Adapun
kedua kurva koordinat lewat sebuah titik berupa \emph{garis
lurus} $\R^3$ (bekukan $u$ atau bekukan $v$: yakni garis pembangun
pelana yang berpembangun ganda), namun $F \neq 0$ di luar sumbunya: jadi garis pembangun
lewat titik generik tak ortogonal. Kedua garis pembangunnya terletak
di \omterm{def:b2:surfaces:tangent}{bidang singgungnya}, yang mereka rentang --- jadi \omterm{def:b2:surfaces:tangent}{bidang
singgungnya} memotong permukaannya sepanjang dua garis penuh, yakni lawan
ekstrem bolanya, yang \omterm{def:b2:surfaces:tangent}{bidang singgungnya} menyentuh di satu
titik belaka. Adapun tanda ``sentuhan orde dua'' antara
sebuah permukaan dan \omterm{def:b2:surfaces:tangent}{bidang singgungnya} merupakan cerita \omterm{thm:b2:curves:frenet2d}{kelengkungan},
yang diangkat pada jilid Tahun ke-3.
\end{example}

\begin{proposition}[Panjang kurva yang digambar pada sebuah
permukaan]\label{prop:b2:surfaces:curvelength}
Bila $\gamma(t) = \sigma(u(t), v(t))$ dengan $t \in [a, b]$ bersifat
$\mathcal{C}^1$, maka
\[
L(\gamma) = \int_a^b
\sqrt{E\,u'^2 + 2F\,u'v' + G\,v'^2}\;\dd t ,
\]
dengan $E, F, G$ dinilai di $(u(t), v(t))$.
\end{proposition}

\begin{proof}
Berlaku $\gamma' = u'\sigma_u + v'\sigma_v$ menurut aturan rantai, jadi
$\norm{\gamma'}^2 = I(u', v') = Eu'^2 + 2Fu'v' + Gv'^2$; lalu integralkan
$\norm{\gamma'}$ (\cref{def:b2:curves:length}).
\end{proof}

\begin{example}[Mengapa pesawat terbang melintasi kutub]\label{ex:b2:surfaces:polarroute}
Dua bandara terletak pada lintang $\varphi_0$ dan bujur
yang berlawanan: $A = \sigma(0, \varphi_0)$ dan $B = \sigma(\pi,
\varphi_0)$ pada bola berjari-jari $R$. Sepanjang paralelnya
(dengan $\varphi \equiv \varphi_0$), \omterm{def:b2:curves:length}{panjangnya} adalah
$\int_0^\pi R\cos\varphi_0\,\dd\theta = \pi R\cos\varphi_0$.
Sedangkan sepanjang rute lewat kutubnya (naik di meridian $\theta = 0$,
lalu turun di meridian $\theta = \pi$), \omterm{def:b2:curves:length}{panjangnya} $2R(\frac\pi2 -
\varphi_0)$. Pada lintang $\varphi_0 = \frac\pi3$ (yakni enam puluh
derajat): rute paralelnya $\pi R/2 \approx 1.571\,R$, sedangkan rute
kutubnya $\pi R/3 \approx 1.047\,R$ --- jadi sepertiga lebih pendek. Bahkan
$\pi\cos\varphi_0 \geq \pi - 2\varphi_0$ pada
$\intcc0{\pi/2}$ (sebab fungsi $\pi\cos\varphi - \pi +
2\varphi$ lenyap di kedua ujungnya dan turunannya $2 -
\pi\sin\varphi$ berganti tanda sekali, jadi ia mula-mula naik
lalu turun, sehingga taknegatif): jadi rute kutubnya tak pernah
kalah. \omterm{def:b2:surfaces:fff}{Bentuk fundamental pertamanya} mengubah pertanyaan navigasi
menjadi dua integral satu baris; dan \cref{exo:b2:surfaces:6} mendorong
gagasannya sampai ke bukti keminimalan yang sejati bagi meridiannya.
\end{example}

\begin{lemma}[Kesamaan Lagrange]\label{lem:b2:surfaces:lagrange}
Untuk setiap $a, b \in \R^3$:
$\norm{a \wedge b}^2 = \norm a^2 \norm b^2 - \langle a, b\rangle^2$.
Khususnya
\[
\norm{\sigma_u \wedge \sigma_v} = \sqrt{EG - F^2}.
\]
\end{lemma}

\begin{proof}
Kedua ruasnya tak berubah bila kita mengganti $b$ dengan komponennya $b_\perp
= b - \frac{\langle a, b\rangle}{\norm a^2}a$ yang ortogonal terhadap $a$
(untuk $a \neq 0$; sedangkan kasus $a = 0$ trivial): ruas kirinya
karena $a \wedge a = 0$, dan ruas kanannya dengan menguraikan $\norm
{b_\perp}^2 = \norm b^2 - \frac{\langle a,b\rangle^2}{\norm a^2}$
beserta $\langle a, b_\perp\rangle = 0$. Jadi cukuplah membuktikan
kesamaannya untuk $a, b$ yang ortogonal, tempat ia berbunyi $\norm{a \wedge
b} = \norm a \norm b$: yang benar, sebab untuk vektor yang ortogonal
hasil kali silangnya bernorma $\norm a\norm b\,\abs{\sin\frac\pi2}$. Adapun
rumus pada displainya adalah kasus $a = \sigma_u$ dan $b = \sigma_v$.
\end{proof}

\begin{remark}
Kesamaan Lagrange mengatakan bahwa $EG - F^2 = \det\begin{psmallmatrix}
E & F\\ F & G\end{psmallmatrix}$ merupakan \emph{\omterm{def:b2:linalg:det}{determinan}
Gram} bagi $(\sigma_u, \sigma_v)$: yakni kuadrat \omterm{def:b2:surfaces:area}{luas}
jajaran genjang yang mereka rentang. Jadi keregularannya, kedefinitan
positif \omterm{def:b2:surfaces:fff}{bentuk fundamental pertamanya}, dan kepositifan
\omterm{def:b2:linalg:det}{determinan} Gramnya adalah tiga ungkapan bagi satu syarat ---
dan itulah sebabnya integran \omterm{def:b2:surfaces:area}{luas} di bawah tak pernah lenyap pada
bagan yang regular.
\end{remark}

\begin{definition}[Luas]\label{def:b2:surfaces:area}
Misalkan $\sigma \colon U \to \R^3$ permukaan $\mathcal{C}^1$ yang regular dan
injektif, serta $K \subseteq U$ daerah \omterm{def:b2:metric:compact}{kompak} yang di atasnya
integral lipat bermakna (\cref{ch:b2:multint}). Adapun
\emph{luas}\index{luas!sebuah permukaan} keping $\sigma(K)$ adalah
\[
\mathcal{A}
= \iint_K \norm{\sigma_u \wedge \sigma_v}\,\dd u\,\dd v
= \iint_K \sqrt{EG - F^2}\;\dd u\,\dd v .
\]
\end{definition}

\begin{remark}[Mengapa rumus ini]
Persegi \omterm{def:b2:curves:length}{panjang} $[u, u + \dd u] \times [v, v + \dd v]$ terpetakan, hingga
orde pertama, ke jajaran genjang yang direntang $\sigma_u\,\dd u$
dan $\sigma_v\,\dd v$, yang berluas $\norm{\sigma_u \wedge
\sigma_v}\,\dd u\,\dd v$: jadi definisinya mengintegralkan faktor
pemuaian \omterm{def:b2:surfaces:area}{luas} lokalnya, persis seperti \omterm{def:b2:curves:length}{panjang busur} yang mengintegralkan
kelajuan lokalnya. Adapun kekonsistenannya dengan \omterm{def:b2:curves:reparam}{penggantian parameter} adalah
\cref{exo:b2:surfaces:8}; sedangkan kekonsistenannya dengan rumus penggantian
peubah bagi integral lipat dibahas pada
\cref{ch:b2:multint}.
\end{remark}

\begin{example}[Dua grafik berbeda, satu
luas]\label{ex:b2:surfaces:samearea}
Di atas cakram satuannya, bandingkanlah mangkuk $z = \frac12(x^2 + y^2)$
dan pelana $z = xy$. Adapun integran luasnya
(\cref{exo:b2:surfaces:5}) adalah
\[
\sqrt{1 + x^2 + y^2}
\qquad\text{dan}\qquad
\sqrt{1 + y^2 + x^2} :
\]
yang identik. Jadi kedua permukaannya --- yang satu melengkung ke arah yang sama
di segala arah, yang lain berbentuk pelana --- berluas persis
sama di atas setiap daerahnya, yakni $\frac{2\pi}3(2\sqrt2 - 1)$
di atas cakram satuannya. Sebab unsur luasnya hanya melihat
\emph{\omterm{def:b2:curves:length}{panjang}} gradiennya, bukan penataan
pembengkokannya; jadi membedakan mangkuknya dari pelananya menuntut
data orde dua (yakni struktur tanda yang dipamerkan pada
\cref{fig:b2:surfaces:saddle}), yang tak terdeteksi oleh pengukuran
\omterm{def:b2:surfaces:area}{luas} seberapa pun banyaknya. Jadi \omterm{def:b2:surfaces:fff}{bentuk fundamental pertama} itu metrik dan buta
terhadap bentuknya; sedangkan bentuk kedua yang melihat bentuknya adalah milik Tahun ke-3.
\end{example}

\begin{example}[Luas bolanya]\label{ex:b2:surfaces:spherearea}
Untuk bagan bola pada \cref{ex:b2:surfaces:standard}:
\[
\sigma_\theta = R(-\sin\theta\cos\varphi,\
\cos\theta\cos\varphi,\ 0),
\qquad
\sigma_\varphi = R(-\cos\theta\sin\varphi,\
-\sin\theta\sin\varphi,\ \cos\varphi),
\]
jadi $E = R^2\cos^2\varphi$, $F = 0$, $G = R^2$ dan $\sqrt{EG - F^2}
= R^2\cos\varphi$. Karena itu
\[
\mathcal{A}
= \int_{-\pi/2}^{\pi/2}\!\!\int_0^{2\pi}
R^2\cos\varphi\;\dd\theta\,\dd\varphi
= 2\pi R^2\,\bigl[\sin\varphi\bigr]_{-\pi/2}^{\pi/2}
= \boxed{4\pi R^2} .
\]
\end{example}

\begin{example}[Kerucutnya, diperiksa terhadap
rumus sekolah]\label{ex:b2:surfaces:conearea}
Untuk kerucut $z = \sqrt{x^2 + y^2}$ di atas anulus $a \leq
\rho \leq b$, rumus grafik pada
\cref{exo:b2:surfaces:5} memberi $1 + f_x^2 + f_y^2 = 2$
(sebab $f_x = x/\rho$ dan $f_y = y/\rho$), jadi
\[
\mathcal A = \sqrt2\,\pi\,(b^2 - a^2) .
\]
Adapun pemeriksaan kekonsistenannya dengan rumus tinggi miring $\pi\rho\ell$
pada \cref{ex:b2:surfaces:revolution}: kerucut penuh yang berjari-jari
alas $b$ dan $a$ punya \omterm{def:b2:surfaces:area}{luas} selimut $\pi b\cdot b\sqrt2$ dan
$\pi a\cdot a\sqrt2$, yang selisihnya persis
$\sqrt2\pi(b^2 - a^2)$. Jadi dua bagan, dua rumus, satu \omterm{def:b2:surfaces:area}{luas}
--- yakni keinvarianan yang dibuktikan pada \cref{exo:b2:surfaces:8}, yang terlihat
di alam liar.
\end{example}

\begin{remark}[Pemeriksaan kewarasan bagi luas]
Tiga pemeriksaan seketika menangkap sebagian besar kesalahan pada perhitungan
\omterm{def:b2:surfaces:area}{luas}. \emph{Penskalaan}: memuaikan sebuah permukaan sebesar $\lambda$
mengalikan $E, F, G$ dengan $\lambda^2$ dan luasnya dengan
$\lambda^2$ --- jadi jawaban yang kebergantungannya pada $R$ tak
kuadratik (seperti $4\pi R^2$) pastilah salah. \emph{Kepositifan
unsurnya}: $\sqrt{EG - F^2}$ mestilah positif tegas pada
interior bagannya; jadi nilai yang lenyap menandai kemerosotan bagannya,
yang harus dipangkas seperti pada kutub bolanya. \emph{Kesimetrian}: sebuah
perhitungan atas keping yang simetris mestilah konsisten dengan
menjumlahkan bagian kongruennya --- jadi setengah bolanya sebaiknya
memberi $2\pi R^2$.
\end{remark}

\begin{example}[Permukaan putaran]\label{ex:b2:surfaces:revolution}
Putarlah kurva $z \mapsto (r(z), 0, z)$ dengan $r > 0$ berkelas
$\mathcal{C}^1$, mengelilingi sumbu-$z$:
\[
\sigma(\theta, z) = (r(z)\cos\theta,\ r(z)\sin\theta,\ z).
\]
Maka $E = r(z)^2$, $F = 0$, $G = 1 + r'(z)^2$, jadi
\[
\mathcal{A} = \int_a^b\!\!\int_0^{2\pi}
r(z)\sqrt{1 + r'(z)^2}\;\dd\theta\,\dd z
= 2\pi\int_a^b r(z)\sqrt{1 + r'(z)^2}\,\dd z ,
\]
yakni rumus klasiknya (yaitu keliling $2\pi r$ dikali unsur \omterm{def:b2:curves:length}{panjang}
miringnya). Untuk kerucut $r(z) = kz$ dengan $z \in [0, h]$: berlaku $\mathcal
A = 2\pi k\sqrt{1 + k^2}\,\frac{h^2}{2} = \pi \rho \ell$ dengan
$\rho = kh$ menyatakan jari-jari alasnya dan $\ell = h\sqrt{1 + k^2}$ tinggi
miringnya --- yakni rumus sekolahnya, yang kini diturunkan alih-alih diterima saja.
\end{example}

\begin{example}[Katenoidnya]\label{ex:b2:surfaces:catenoid}
Putarlah katenari $r(z) = \cosh z$ dengan $z \in \intcc{-1}{1}$
mengelilingi sumbunya: maka \emph{katenoid} hasilnya punya, menurut
rumus putarannya dan $1 + \sinh^2 z = \cosh^2 z$,
\[
\mathcal A = 2\pi\int_{-1}^1\cosh z\,\sqrt{1 + \sinh^2z}\,
\dd z = 2\pi\int_{-1}^1\cosh^2z\,\dd z
= \pi\bigl[z + \sinh z\cosh z\bigr]_{-1}^{1},
\]
yakni
\[
\mathcal A = 2\pi + \pi\sinh 2 \approx 17.68 .
\]
Faktor miringnya $\sqrt{1 + r'^2}$ melebur dengan jari-jarinya
menjadi kuadrat sempurna --- yakni kesamaan yang sama yang membuat
\omterm{def:b2:curves:length}{panjang busur} katenari menjadi dasar pada bab kurva. Ini
bukan kebetulan aljabar: sebab di antara semua permukaan putaran
yang merentang kedua lingkaran perbatasannya, katenoidnya meminimumkan
luasnya (yakni bentuk selaput sabun di antara dua cincin), dan
sifat variasional inilah yang persis menyendirikan
$\cosh$; adapun jilid Tahun ke-3 membuktikannya lewat kalkulus
variasi.
\end{example}

\begin{figure}[ht]
\centering
\begin{tikzpicture}[scale=1.05]
  % A saddle surface z = x^2 - y^2 drawn as a wire mesh in oblique projection
  % projection: (x,y,z) -> (x + 0.45*y, z + 0.35*y)
  \begin{scope}
    % u-lines (x varies, y fixed)
    \foreach \yy in {-1,-0.5,0,0.5,1} {
      \draw[ocDarkBlue!70, thick, domain=-1:1, smooth, samples=25]
        plot ({\x + 0.45*\yy}, {\x*\x - \yy*\yy + 0.35*\yy});
    }
    % v-lines (y varies, x fixed)
    \foreach \xx in {-1,-0.5,0,0.5,1} {
      \draw[ocMint!80!black, thick, domain=-1:1, smooth, samples=25]
        plot ({\xx + 0.45*\x}, {\xx*\xx - \x*\x + 0.35*\x});
    }
    % tangent plane at origin: z = 0 -> parallelogram
    \draw[ocRed, dashed]
      (-0.8 - 0.36, -0.28) -- (0.8 - 0.36, -0.28)
      -- (0.8 + 0.36, 0.28) -- (-0.8 + 0.36, 0.28) -- cycle;
    \fill (0,0) circle (0.045) node[below right, font=\small] {$M_0$};
    % normal vector at origin
    \draw[->, thick] (0,0) -- (0,0.9)
      node[right, font=\small] {$n$};
  \end{scope}
\end{tikzpicture}
\caption{Pelana $z = x^2 - y^2$ di dekat titik asalnya, beserta
kurva koordinatnya (yakni kurva-$u$ yang biru dan kurva-$v$ yang hijau),
\omterm{def:b2:surfaces:tangent}{bidang singgungnya} di $M_0 = (0,0,0)$ (yang bergaris putus-putus) dan \omterm{def:b2:surfaces:tangent}{normal satuannya} $n$.
Permukaannya menyeberangi \omterm{def:b2:surfaces:tangent}{bidang singgungnya} --- yakni analogi
dua dimensi bagi sebuah belokan.}
\label{fig:b2:surfaces:saddle}
\end{figure}

\begin{remark}[Jebakan yang sering muncul]
(i) \emph{Kesingularan bagan bukan kesingularan permukaan}:
sebab bagan bolanya merosot di kutubnya
($\cos\varphi = 0$), tetapi bolanya mulus sempurna di sana
--- sebab bagan lain (dengan peran sumbunya dipertukarkan) bersifat regular
di kutubnya. Jadi sebelum mengumumkan sebuah titik singular, cobalah
parameterisasi kedua. (ii) \emph{Keregularan $\sigma$ menyangkut
parameterisasinya, bukan petanya}: sebab $\sigma(u, v) = (u^3, v,
0)$ tak regular sepanjang $u = 0$ walaupun petanya berupa sebuah
bidang. (iii) Adapun rumus luasnya menuntut $\sigma$ yang
\emph{injektif} pada $K$: sebab bagan yang meliputi sebuah keping dua kali
mencacahnya dua kali (yakni $\theta$ yang berjalan atas $\intcc0{4\pi}$
melipatduakan \omterm{def:b2:surfaces:area}{luas} bolanya). (iv) Adapun \omterm{def:b2:surfaces:tangent}{normal satuannya} terdefinisi
hingga tandanya oleh permukaannya tetapi \emph{dipilih} oleh bagannya
(yakni oleh urutan $u, v$); jadi pernyataan yang melibatkan orientasi harus menetapkan
pilihan itu. (v) Akhirnya, $EG - F^2 > 0$ bukanlah hipotesis
tambahan: sebab ia persis keregularannya, menurut kesamaan
Lagrange --- jadi bila ia lenyap di suatu tempat, masalahnya ada pada
bagannya, dan tak ada rumus \omterm{def:b2:surfaces:area}{luas} atau \omterm{def:b2:surfaces:tangent}{bidang singgung} yang berlaku di sana.
\end{remark}

\begin{remark}[Pandangan ke depan di dalam jilid ini]
Pautan ke depan dari sini. Unsur luasnya $\sqrt{EG -
F^2}\,\dd u\,\dd v$ adalah Jacobi dua dimensi yang
menyamar, dan \cref{ch:b2:multint} membuat analoginya eksak
lewat teorema penggantian peubah --- jadi integral
permukaan di sana adalah \omterm{def:b2:surfaces:area}{luas} bab ini dengan sebuah integran di
atasnya. Adapun \omterm{def:b2:surfaces:fff}{bentuk fundamental pertamanya} adalah medan \omterm{def:b2:quadratic:def}{bentuk
kuadratik} yang positif, yang ditangani \omterm{def:b2:funcseq:def}{titik demi titik} lewat perkakas
\cref{ch:b2:quadratic}, dan teorema spektralnya juga menggerakkan
penggolongan \omterm{pb:b2:surfaces:1}{kuadrik} pada soal akhir pekan bab ini. Lalu garis
normalnya menggerakkan masalah ekstremum pada himpunan kendala
(\cref{ex:b2:surfaces:closest}), yakni benih geometris bagi
metode pengganda Lagrange yang disketsakan bersama
\cref{thm:b2:diffcalc:extrema}.
\end{remark}

\begin{omfigure}
\begin{tikzpicture}[scale=0.8]
  \draw[omDef, thick] (0,0) ellipse (1.5 and 0.95);
  \draw[omThm] (0,0) ellipse (1.5 and 0.35);
  \node[below, font=\small] at (0,-1.25) {elipsoid};
\end{tikzpicture}
\qquad
\begin{tikzpicture}[scale=0.8]
  \draw[omDef, thick] plot[domain=-1.1:1.1, samples=30]
    ({0.6*(exp(\x)+exp(-\x))/2}, {\x});
  \draw[omDef, thick] plot[domain=-1.1:1.1, samples=30]
    ({-0.6*(exp(\x)+exp(-\x))/2}, {\x});
  \draw[omThm] (0,0) ellipse (0.6 and 0.16);
  \draw[omDef] (0,1.1) ellipse (1.0 and 0.22);
  \draw[omDef] (0,-1.1) ellipse (1.0 and 0.22);
  \node[below, font=\small] at (0,-1.55)
    {hiperboloid berdaun satu};
\end{tikzpicture}
\qquad
\begin{tikzpicture}[scale=0.8]
  \draw[omDef, thick] plot[domain=-1.05:1.05, samples=30]
    ({\x}, {1.1*\x*\x - 1.1});
  \draw[omThm] (0,0.055) ellipse (1.0 and 0.22);
  \node[below, font=\small] at (0,-1.45)
    {paraboloid eliptik};
\end{tikzpicture}
\qquad
\begin{tikzpicture}[scale=0.8]
  \draw[omDef, thick] plot[domain=-1:1, samples=30]
    ({\x}, {0.55*\x*\x});
  \draw[omThm, thick] plot[domain=-1:1, samples=30]
    ({0.45*\x}, {-0.55*\x*\x + 0.18*\x});
  \node[below, font=\small] at (0,-1.15)
    {paraboloid hiperbolik};
\end{tikzpicture}

{\small Empat dari sembilan permukaan \omterm{pb:b2:surfaces:1}{kuadrik} yang digolongkan pada
soal akhir pekan, yang disketsakan lewat siluetnya dan sebuah kurva
aras (merah): yakni elipsoid yang terbatas, hiperboloid berdaun
satu yang berpembangun ganda beserta pinggangnya, mangkuk
paraboloid eliptiknya, dan pelananya, yang kedua irisan parabolanya
membengkok ke arah berlawanan.}
\end{omfigure}

\begin{remark}[Di mana ini dipakai]
Unsur luasnya $\norm{\sigma_u\wedge\sigma_v}\,\dd u\,\dd v$
adalah ukuran integral permukaan pada \cref{ch:b2:multint}, tempat
ia bertemu rumus Green; sedangkan \omterm{def:b2:surfaces:fff}{bentuk fundamental pertamanya} adalah
prototipe sebuah medan \omterm{def:b2:quadratic:def}{bentuk kuadratik}, yang dikaji \omterm{def:b2:funcseq:def}{titik demi titik}
lewat perkakas \cref{ch:b2:quadratic}; dan soal akhir pekan bab
ini menggolongkan semua permukaan \omterm{pb:b2:surfaces:1}{kuadrik} lewat teorema
spektralnya. Adapun jilid Tahun ke-3 kembali ke permukaan lewat
bentuk \omterm{def:b2:diffcalc:differential}{diferensial} dan teorema divergensinya, sedangkan \omterm{thm:b2:curves:frenet2d}{kelengkungan}
hakikinya --- yakni apa yang diketahui $E, F, G$ tentang pembengkokannya --- adalah
gerbang menuju geometri \omterm{def:b2:diffcalc:differential}{diferensial} yang sesungguhnya.
\end{remark}

\section{Latihan}

\begin{exercise}[$\star$]\label{exo:b2:surfaces:1}
Tunjukkan bahwa \omterm{def:b2:surfaces:tangent}{bidang singgung} kerucut $z = \sqrt{x^2 + y^2}$
(dikurangi puncaknya) semuanya melalui puncaknya. \emph{(Parameterkan lewat
$\sigma(\theta, r) = (r\cos\theta, r\sin\theta, r)$ dengan $r > 0$.)}
\end{exercise}

\begin{exercise}[$\star$]\label{exo:b2:surfaces:2}
Carilah \omterm{def:b2:surfaces:tangent}{bidang singgung} elipsoid $\frac{x^2}{a^2} +
\frac{y^2}{b^2} + \frac{z^2}{c^2} = 1$ di sebuah titik $(x_0, y_0,
z_0)$ pada permukaannya.
\end{exercise}

\begin{exercise}[$\star$]\label{exo:b2:surfaces:3}
Hitunglah $E, F, G$ bagi helikoid $\sigma(u, v) = (v\cos u,\
v\sin u,\ au)$ dengan $a > 0$, beserta \omterm{def:b2:surfaces:area}{luas} keping $0 \leq u \leq
2\pi$, $0 \leq v \leq 1$, sebagai sebuah integral (lalu nilailah ia dengan memakai
$\int\sqrt{v^2 + a^2}\,\dd v = \frac12\bigl(v\sqrt{v^2+a^2} +
a^2\ln(v + \sqrt{v^2 + a^2})\bigr) + C$).
\end{exercise}

\begin{exercise}[$\star\star$]\label{exo:b2:surfaces:4}
(Torus) Parameterkan torus yang diperoleh dengan memutar lingkaran
berpusat $(R, 0, 0)$ dan berjari-jari $r < R$ pada bidang-$xz$ mengelilingi
sumbu-$z$:
\[
\sigma(\theta, \psi) = \bigl((R + r\cos\psi)\cos\theta,\
(R + r\cos\psi)\sin\theta,\ r\sin\psi\bigr).
\]
Hitunglah $E, F, G$, periksalah keregularannya, lalu tunjukkan bahwa luasnya
$4\pi^2 R r$ (yakni Pappus: keliling rata-ratanya $2\pi R$ dikali \omterm{def:b2:curves:length}{panjang}
lingkarannya $2\pi r$).
\end{exercise}

\begin{exercise}[$\star\star$]\label{exo:b2:surfaces:5}
Tunjukkan bahwa \omterm{def:b2:surfaces:area}{luas} grafik $f \in \mathcal{C}^1(K)$ adalah
$\iint_K \sqrt{1 + f_x^2 + f_y^2}\,\dd x\,\dd y$, lalu hitunglah ia
bagi keping paraboloid $z = \frac12(x^2 + y^2)$ di atas cakram
$x^2 + y^2 \leq 1$ (lewat koordinat kutub,
\cref{ch:b2:multint}).
\end{exercise}

\begin{exercise}[$\star\star$]\label{exo:b2:surfaces:6}
Sebuah kurva tergambar $\gamma(t) = \sigma(u(t), v(t))$ pada bola
berjari-jari $R$ (dengan bagan bolanya) punya $u = \theta(t)$ dan $v =
\varphi(t)$. Tulislah \omterm{def:b2:curves:length}{panjangnya} sebagai integral dalam $\theta, \varphi$
lalu buktikan bahwa di antara kurva yang menghubungkan dua titik pada meridian
yang sama $\theta = \theta_0$, busur meridiannya yang terpendek.
\emph{(Batasilah integrannya dari bawah oleh $R\abs{\varphi'}$.)}
\end{exercise}

\begin{exercise}[$\star\star\star$]\label{exo:b2:surfaces:7}
(Garis normal sebuah bola) Misalkan $S$ permukaan aras yang regular
$\{f = c\}$, \omterm{def:b2:metric:connected}{terhubung}, dan semua garis normalnya melalui sebuah
titik tetap $\Omega$. Tunjukkan bahwa $S$ termuat pada sebuah bola
yang berpusat $\Omega$. \emph{(Tunjukkan bahwa $\norm{M - \Omega}^2$ punya
turunan nol sepanjang setiap kurva yang digambar pada $S$.)}
\end{exercise}

\begin{exercise}[$\star\star\star$]\label{exo:b2:surfaces:8}
(Luasnya geometris) Misalkan $\Phi \colon U' \to U$ difeomorfisma
$\mathcal{C}^1$ antara himpunan \omterm{def:b2:metric:topology}{terbuka} $\R^2$ dan $\tilde\sigma =
\sigma \circ \Phi$. Tunjukkan bahwa
\[
\norm{\tilde\sigma_{u'} \wedge \tilde\sigma_{v'}}
= \abs{\det J_\Phi}\,
\norm{(\sigma_u \wedge \sigma_v)\circ\Phi} ,
\]
lalu turunkan, dengan memakai rumus penggantian peubah pada
\cref{ch:b2:multint}, bahwa \omterm{def:b2:surfaces:area}{luas} pada
\cref{def:b2:surfaces:area} tak bergantung pada parameterisasi regular
yang dipilih.
\end{exercise}

\begin{exercise}[$\star$]\label{exo:b2:surfaces:9}
(Teorema kotak-topi Archimedes) Pada bola berjari-jari $R$,
\emph{zona} antara lintang yang $z_1 \leq z \leq z_2$
(dengan $-R \leq z_1 < z_2 \leq R$) berluas $2\pi R\,(z_2 - z_1)$:
buktikanlah lewat bagan bolanya, lalu simpulkan bahwa \omterm{def:b2:surfaces:area}{luas} sebuah zona
hanya bergantung pada tingginya --- jadi mengiris sebuah jeruk menjadi
irisan setebal sama memberi jumlah kulit yang sama.
\end{exercise}

\begin{exercise}[$\star\star$]\label{exo:b2:surfaces:10}
Tunjukkan bahwa setiap garis normal permukaan putaran
$\sigma(\theta, z) = (r(z)\cos\theta,\ r(z)\sin\theta,\ z)$
(dengan $r > 0$ berkelas $\mathcal C^1$) bertemu sumbu
putarannya, lalu tentukan letak titik potongnya.
\end{exercise}

\begin{exercise}[$\star\star$]\label{exo:b2:surfaces:11}
(Membuka gulungan silinder) Bagan $\sigma(u, v) = (\cos u,\
\sin u,\ v)$ bagi silinder satuannya punya $E = G = 1$ dan $F = 0$:
periksalah ini, lalu jelaskan mengapa setiap kurva tergambar $t \mapsto
\sigma(u(t), v(t))$ berpanjang sama dengan kurva bidang $t
\mapsto (u(t), v(t))$. Turunkan bahwa heliks dari $(1, 0, 0)$
ke $(1, 0, 2\pi c)$ yang membuat satu putaran berpanjang
$2\pi\sqrt{1 + c^2}$, dan bahwa tak ada kurva tergambar dengan
titik ujung yang sama dan satu putaran penuh yang lebih pendek.
\end{exercise}

\begin{exercise}[$\star\star\star$]\label{exo:b2:surfaces:12}
Misalkan $S = \{f = c\}$ permukaan aras regular yang \omterm{def:b2:metric:compact}{kompak} dan
$M_0 \in S$ sebuah titik yang berjarak maksimum dari titik asalnya. Tunjukkan
bahwa $\nabla f(M_0)$ segaris dengan $\vect{OM_0}$ --- jadi
normal di titik terjauhnya bersifat radial. Terapkanlah pada elipsoid
$\frac{x^2}{a^2} + \frac{y^2}{b^2} + \frac{z^2}{c^2} = 1$ (dengan $a >
b > c > 0$): carilah semua titik yang normalnya radial, lalu
kenalilah yang terjauhnya.
\end{exercise}

\section{Soal: penggolongan kuadrik pada
$\R^3$}

\begin{problem}[{Soal akhir pekan --- setiap permukaan kuadrik,
disortir lewat teorema spektral}]\label{pb:b2:surfaces:1}
Sebuah \emph{kuadrik}\index{kuadrik} adalah himpunan nol di $\R^3$ bagi
sebuah polinomial berderajat dua
\[
q(X) = X^{\mathsf T}\!AX + 2\,\langle b, X\rangle + c,
\qquad A \in \mathcal S_3(\R),\ A \neq 0,\ b \in \R^3,\
c \in \R .
\]
Adapun permukaan pada gambar bab ini --- yakni bola,
elipsoid, pelana, kerucut, silinder --- semuanya \omterm{pb:b2:surfaces:1}{kuadrik}.
Soal ini menggolongkan \emph{semuanya}: sebab teorema spektral
(\cref{thm:b2:quadratic:spectral}) meluruskan bagian
kuadratiknya, penggeseran \omterm{def:b2:affine:subspace}{afin} (\cref{ch:b2:affine}) menyerap
bagian linearnya, dan yang tersisa adalah daftar bentuk normal yang
pendek dan \omterm{def:b2:metric:complete}{lengkap}.

\medskip
\noindent\textbf{Bagian I --- Mesin reduksinya.}
\begin{enumerate}
  \item Misalkan $X = PY + t$ dengan $P \in O(3)$ dan $t \in \R^3$ (yakni
        penggantian koordinat yang kaku). Tunjukkan bahwa $q(PY + t) =
        Y^{\mathsf T}\!A'Y + 2\langle b', Y\rangle + c'$ dengan
        \[
        A' = P^{\mathsf T}\!AP, \qquad
        b' = P^{\mathsf T}(At + b), \qquad
        c' = q(t) .
        \]
        Lalu turunkan bahwa \omterm{def:b2:reduction:eigen}{spektrum} $A$ (sehingga rank dan
        signaturnya) merupakan invarian kaku bagi persamaannya, dan
        jelaskanlah mengapa persamaan sebuah \omterm{pb:b2:surfaces:1}{kuadrik} hanya
        tertentukan hingga sebuah faktor skalar tak nol.
  \item Dengan memakai teorema spektralnya, tunjukkanlah bahwa setelah sebuah perputaran
        persamaannya menjadi $\sum_i \lambda_i y_i^2 +
        2\sum_i\beta_iy_i + c = 0$ dengan $\lambda_1, \lambda_2,
        \lambda_3$ menyatakan \omterm{def:b2:reduction:eigen}{nilai eigen} $A$.
  \item Untuk setiap $i$ yang $\lambda_i \neq 0$, serap
        $\beta_iy_i$ lewat sebuah penggeseran ($y_i \mapsto y_i -
        \beta_i/\lambda_i$). Tulislah persamaan tereduksinya ketika
        $\operatorname{rank} A = r$: $\sum_{i\leq r}\lambda_i
        z_i^2 + 2\sum_{i > r}\beta_iz_i + c'' = 0$.
  \item Sebuah \emph{pusat} \omterm{pb:b2:surfaces:1}{kuadrik} berpersamaan $q = 0$ adalah
        titik $\Omega$ dengan $q(2\Omega - X) = q(X)$ untuk setiap
        $X$: jadi pencerminan titik pada $\Omega$ memelihara
        persamaannya, sehingga permukaannya. Tunjukkan bahwa $q(2\Omega - X)
        - q(X) = -4\langle A\Omega + b, X\rangle + 4\langle
        A\Omega + b, \Omega\rangle$, lalu turunkan: bahwa pusatnya
        persis merupakan penyelesaian $A\Omega = -b$; jadi ia
        ada bila dan hanya bila $b \in \operatorname{im}A$, dan pusatnya
        tunggal bila dan hanya bila $A$ terbalikkan.
\end{enumerate}

\medskip
\noindent\textbf{Bagian II --- \omterm{pb:b2:surfaces:1}{Kuadrik} berpusat
(dengan $\operatorname{rank}A = 3$).} Di sini persamaan tereduksinya adalah
$\lambda_1z_1^2 + \lambda_2z_2^2 + \lambda_3z_3^2 = \delta$.
\begin{enumerate}[resume]
  \item Dengan mengalikan $-1$ bila perlu, andaikan sekurangnya dua
        $\lambda_i > 0$. Cacahlah kemungkinannya:
        yakni signatur $(3, 0)$ dengan $\delta > 0$, $= 0$, $< 0$,
        dan signatur $(2, 1)$ dengan $\delta > 0$, $= 0$, $<
        0$; lalu namailah keenam himpunan hasilnya (elipsoid, titik,
        himpunan kosong, hiperboloid berdaun satu, kerucut, hiperboloid
        berdaun dua) kemudian tempatkanlah masing-masingnya dalam bentuk
        normal Euklidesnya ($\frac{x^2}{a^2} + \frac{y^2}{b^2} +
        \frac{z^2}{c^2} = 1$, dan seterusnya).
  \item Golongkanlah $x^2 + y^2 + z^2 + 4xy + 4yz + 4zx = 1$:
        tunjukkan $A = 2J - I$ dengan $J$ matriks serba-satu,
        hitunglah \omterm{def:b2:reduction:eigen}{spektrumnya} $\{5, -1, -1\}$, lalu kenalilah sebuah
        hiperboloid putaran berdaun dua terhadap sumbu
        $\R(1,1,1)$.
  \item (Garis pembangunnya) Untuk hiperboloid berdaun satu
        $\frac{x^2}{a^2} + \frac{y^2}{b^2} - \frac{z^2}{c^2}
        = 1$, faktorkanlah
        \[
        \Bigl(\frac xa - \frac zc\Bigr)
        \Bigl(\frac xa + \frac zc\Bigr)
        = \Bigl(1 - \frac yb\Bigr)\Bigl(1 + \frac yb\Bigr)
        \]
        lalu hasilkanlah dua keluarga berparameter satu berisi garis
        lurus yang terletak pada permukaannya.
  \item Tunjukkan bahwa lewat setiap titik hiperboloid berdaun
        satu melintas tepat satu garis dari tiap keluarganya:
        jadi permukaannya \emph{berpembangun ganda}.
  \item Adapun \emph{kerucut asimtotik} bagi hiperboloid berdaun
        satu adalah $C : \frac{x^2}{a^2} + \frac{y^2}{b^2}
        - \frac{z^2}{c^2} = 0$. Dengan $\rho^2 = \frac{x^2}{a^2}
        + \frac{y^2}{b^2}$, tunjukkanlah bahwa sembarang titik
        hiperboloidnya berjarak paling banyak $c\bigl(\rho -
        \sqrt{\rho^2 - 1}\bigr) = \frac{c}{\rho +
        \sqrt{\rho^2 - 1}}$ dari $C$, jadi permukaannya memeluk
        kerucutnya di tak hingga. Lalu apakah irisan
        hiperboloidnya oleh bidang $x = \pm a$?
\end{enumerate}

\medskip
\noindent\textbf{Bagian III --- Rank $2$ dan rank $1$:
paraboloid, silinder, bidang.}
\begin{enumerate}[resume]
  \item Andaikan $\operatorname{rank}A = 2$, katakanlah $\lambda_1,
        \lambda_2 \neq 0 = \lambda_3$. Mulailah dari pertanyaan
        3, lalu pecahlah menjadi dua kasus menurut $\beta_3 \neq 0$
        (yakni tanpa pusat, menurut pertanyaan 4) atau $\beta_3 = 0$ (yakni sebuah garis
        pusat), lalu reduksikanlah menjadi
        \[
        \lambda_1z_1^2 + \lambda_2z_2^2 + 2\beta_3z_3 = 0
        \qquad\text{atau}\qquad
        \lambda_1z_1^2 + \lambda_2z_2^2 + c'' = 0 :
        \]
        yakni paraboloid eliptik/hiperbolik pada kasus pertamanya,
        dan silinder atas konik berpusat (atau sepasang
        bidang yang berpotongan, sebuah garis, atau himpunan kosong) pada
        kasus keduanya.
  \item Tunjukkan bahwa pelana $z = xy$ merupakan paraboloid
        hiperbolik: putarlah sebesar $\pi/4$ pada bidang-$xy$ untuk
        mencapai $z = \tfrac12(u^2 - v^2)$, yakni permukaan pada
        \cref{fig:b2:surfaces:saddle} hingga penskalaan.
  \item Tunjukkan bahwa pelana $z = xy$ mengusung kedua keluarga
        garis $\{x = x_0,\ z = x_0y\}$ dan $\{y = y_0,\ z
        = xy_0\}$, dengan tepat satu garis dari tiap keluarganya lewat
        setiap titiknya: yakni \omterm{pb:b2:surfaces:1}{kuadrik} berpembangun ganda yang kedua.
  \item Golongkanlah $x^2 + y^2 - 2x + 4y + 3 = 0$ di $\R^3$
        (lengkapkan kuadratnya; lalu kenalilah sebuah silinder lingkaran
        tegak, kemudian berikanlah sumbu dan jari-jarinya).
  \item Kini misalkan $\operatorname{rank}A = 1$, katakanlah $\lambda_1
        \neq 0 = \lambda_2 = \lambda_3$. Dengan memutar di dalam
        bidang kernelnya lalu menggeser, reduksikanlah menjadi
        \[
        \lambda_1z_1^2 + 2\beta z_2 = 0 \quad (\beta \neq 0)
        \qquad\text{atau}\qquad
        \lambda_1z_1^2 + c'' = 0 :
        \]
        yakni sebuah silinder parabolik, atau sepasang bidang sejajar, sebuah
        bidang ganda, atau himpunan kosong. Golongkanlah $(x + y)^2 =
        z$ selengkapnya (yakni bentuk normalnya, dan sumbu
        keinvarianan geserannya).
\end{enumerate}

\medskip
\noindent\textbf{Bagian IV --- Teorema penggolongannya.}
\begin{enumerate}[resume]
  \item Rangkailah Bagian I--III menjadi sebuah teorema: \emph{bahwa setiap
        \omterm{pb:b2:surfaces:1}{kuadrik} $\R^3$ dipetakan oleh sebuah gerak kaku ke
        tepat satu bentuk normal}. Daftarkanlah ketujuh belas tipe
        \omterm{def:b2:affine:subspace}{afinnya} (cacahlah varian kosong dan himpunan yang merosot),
        lalu sendirikanlah kesembilan \emph{permukaan} \omterm{pb:b2:surfaces:1}{kuadriknya}:
        yakni elipsoid, hiperboloid berdaun satu dan berdaun dua, kerucut,
        paraboloid eliptik dan hiperbolik, serta silinder eliptik,
        hiperbolik dan parabolik.
  \item Tulislah \emph{algoritma} penggolongannya: yakni diberikan $(A,
        b, c)$, besaran mana yang kauhitung, dalam urutan mana, dan
        cabang mana yang memutuskan tipe mana? Benarkanlah
        bahwa tiap langkahnya efektif (yakni \omterm{def:b2:reduction:eigen}{nilai eigen} sebuah matriks
        \omterm{def:b2:quadratic:adjoint}{simetrik} $3\times3$, rank, dan keterselesaian
        $A\Omega = -b$).
  \item Jalankan algoritmanya pada $x^2 + y^2 + z^2 - 2xy - 2yz -
        2zx = 1$: tunjukkan bahwa $A = 2I - J$ berspektrum $\{2, 2,
        -1\}$ lalu simpulkan: yakni hiperboloid putaran berdaun satu
        terhadap $\R(1,1,1)$.
  \item Jalankan ia pada $x^2 + y^2 - z^2 - 2x + 4y + 2z + 4 = 0$:
        carilah pusatnya lalu kenalilah \omterm{pb:b2:surfaces:1}{kuadriknya}.
  \item Jalankan ia pada $x^2 + 4xy + y^2 = 2z$: diagonalkan
        blok-$xy$-nya ($u = \frac{x+y}{\sqrt2}$, $v =
        \frac{x-y}{\sqrt2}$) lalu kenalilah \omterm{pb:b2:surfaces:1}{kuadriknya}.
  \item Euklides lawan \omterm{def:b2:affine:subspace}{afin}. Tunjukkan bahwa dua \omterm{pb:b2:surfaces:1}{kuadrik} berpusat
        dalam bentuk normalnya setara secara \emph{kaku}
        bila dan hanya bila daftar koefisiennya sama (hingga
        permutasi dan sebuah skalar positif bersama pada
        persamaannya), sedangkan secara \emph{\omterm{def:b2:affine:subspace}{afin}}, hanya
        data signaturnya yang bertahan: jadi setiap elipsoid adalah peta \omterm{def:b2:affine:subspace}{afin}
        bola bundarnya. Lalu teorema mana yang menjamin bahwa
        signaturnya tak dapat berubah sepanjang jalannya
        (\cref{thm:b2:quadratic:sylvester})?
\end{enumerate}

\medskip
\noindent\textbf{Bagian V --- Panennya.}
\begin{enumerate}[resume]
  \item Tunjukkan bahwa setiap irisan sebuah \omterm{pb:b2:surfaces:1}{kuadrik} oleh bidang
        \omterm{def:b2:affine:subspace}{afin} merupakan sebuah konik (yang boleh merosot) pada bidang itu.
        Kenalilah irisan pelana $z = xy$ oleh
        bidang $z = c$ (untuk $c \neq 0$ dan $c = 0$).
  \item Permukaan \omterm{pb:b2:surfaces:1}{kuadrik} mana yang memuat garis lurus? Tunjukkan
        bahwa elipsoid, hiperboloid berdaun dua dan
        paraboloid eliptik tak memuat satu pun \emph{(batasilah $q$
        ke sebuah garis lalu pakailah ketaksamaan Cauchy--Schwarz bagi
        kasus berdaun dua)}; bahwa kerucut dan silinder
        \omterm{def:b2:structures:generated}{dibangun} satu keluarga; lalu simpulkan bahwa permukaan \omterm{pb:b2:surfaces:1}{kuadrik}
        yang berpembangun ganda persis merupakan hiperboloid berdaun satu
        dan paraboloid hiperbolik.
  \item Ketika hanya tipe \emph{\omterm{def:b2:affine:subspace}{afinnya}} yang dikehendaki, \omterm{thm:b2:quadratic:gauss}{reduksi
        Gauss} (\cref{thm:b2:quadratic:gauss}) lebih murah
        daripada mendiagonalkan. Kerjakanlah ulang pertanyaan 17 lewat algoritma
        Gauss lalu periksalah signaturnya $(2, 1)$; dan informasi
        Euklides apa yang hilang oleh Gauss?
  \item Semua \omterm{def:b2:reduction:eigen}{nilai eigen} matriks \omterm{def:b2:quadratic:adjoint}{simetrik} bersifat real; tunjukkanlah
        bahwa karenanya \emph{tanda} \omterm{def:b2:reduction:eigen}{nilai eigen}
        $A$ dapat dibaca dari \omterm{def:b2:reduction:charpoly}{polinomial karakteristiknya}
        lewat aturan tanda Descartes, lalu periksalah pada
        pertanyaan 6: bahwa $\chi_A(\lambda) = \lambda^3 - 3\lambda^2
        - 9\lambda - 5$ punya tepat satu pergantian tanda, sehingga
        bersignatur $(1, 2)$.
  \item Rangkuman. Rangkumlah algoritmanya dalam beberapa baris;
        lalu nyatakanlah peran persis yang dimainkan (i) teorema
        spektral, (ii) persamaan pusatnya $A\Omega = -b$,
        (iii) teorema inersia Sylvester, dan (iv) \omterm{thm:b2:quadratic:gauss}{reduksi
        Gauss}. Lalu apa yang dihasilkan mesin yang sama di
        $\R^2$, dan apa yang berubah di $\R^n$?

\end{enumerate}
\end{problem}
